\subsection{Software Maintenance}
\begin{frame}{Wisdom on Maintenance} % slided copied from se1-04a. reuse?
	\begin{fancycolumns}
		\pic[width=\linewidth,trim=0 360 0 50,clip]{people/karolina-szczur}
		\vspace{-7mm}

		\begin{note}{Karolina Szczur \mysource{\href{https://twitter.com/codewisdom/status/1141686386103332864}{twitter.com}}}
			\mycite{Writing software as if we are the only person that ever has to comprehend it is one of the biggest mistakes and false assumptions that can be made.}
		\end{note}
		\nextcolumn
		\pic[width=\linewidth,trim=450 235 450 120,clip]{changes/psychopath}
		\vspace{-7mm}

		\begin{note}{John F. Woods \mysource{\href{https://groups.google.com/g/comp.lang.c++/c/rYCO5yn4lXw/m/oITtSkZOtoUJ}{google.com}}}
			\mycite{Always code as if the guy who ends up maintaining your code will be a violent psychopath who knows where you live.}
		\end{note}
	\end{fancycolumns}
\end{frame}

\begin{frame}{\insertsubsection\ \mytitlesource{\ludewiglichter}} % slided copied from se1-04a. reuse?
	\begin{fancycolumns}
		\begin{note}{Motivation}
			\begin{itemize}
				\item for software: no compensation of deterioration, repair, spare parts
				\item corrections (especially shortly after first delivery)
				\item modification and reconstruction
			\end{itemize}
		\end{note}
		\uncover<2->{\begin{definition}{Operation and Maintenance Phase \mysource{\ieeesixten}}
				\mycite{The period of time in the software life cycle during which a software product is employed in its operational environment, monitored for satisfactory performance, and modified as necessary to correct problems or to respond to changing requirements.}
			\end{definition}}
		\nextcolumn
		\uncover<3->{\begin{definition}{Maintenance \mysource{\ieeesixten}}
				\mycite{The process of modifying a software system or component after delivery to correct faults, improve performance or other attributes, or adapt to a changed environment.}
			\end{definition}}
		\uncover<3->{\begin{note}{Kinds of Maintenance \mysource{\ludewiglichter}}
				corrective maintanance, perfective maintenance, adaptive maintenance, preventive maintenance
			\end{note}}
	\end{fancycolumns}
\end{frame}

\subsection{Evolution and Maintenance}
\begin{frame}{\insertsubsection\ \mytitlesource{\ludewiglichter}}
	\slideEvolutionAndMaintenance
\end{frame}
% TODO Operations

\subsection{Types of Maintenance}
\begin{frame}{Corrective Maintenance}
	\begin{fancycolumns}
		\begin{definition}{\insertsubsection\mysource{\lientzswanson}}
			\mycite{Maintenance performed to correct faults in software.} \hfill \deutsch{korrektive Wartung}
		\end{definition}
		\begin{example}{Xiaomi SU7 Crash\mysource{\href{https://www.nytimes.com/2025/04/01/business/xiaomi-driverless-electric-car-china.html}{nytimes.com}}}
			\begin{itemize}
				\item car crashed into guardrail with 97 km/h and caught fire
				\item despite advanced driver assistance turned on
				\item investigation may reveal need for corrections
			\end{itemize}
			\centering\pic[width=.7\linewidth]{automotive/xiaomi-su7}
		\end{example}
		\nextcolumn
		\begin{example}{Ariane 5 Software Failure\mysource{\href{https://doi.ieeecomputersociety.org/10.1109/2.562936}{J{\'e}z{\'e}quel and Meyer, 1997}}}
			\begin{itemize}
				\item explosion about 40 seconds after takeoff
				\item software component of Ariane 4 was reused
				\item value exceeding 16-bit signed integer
				\item unnecessary component removed for launch one year later with new hardware
			\end{itemize}
			\centering\pic[width=.85\linewidth]{failures/ariane5-explosion} {\tiny\copyright{} ESA} % TODO how to add copyright? ESA
		\end{example}
	\end{fancycolumns}
\end{frame}

\xkcdframe{1172} % No more CPU heating as breaking change

\begin{frame}{Perfective Maintenance}
	\begin{fancycolumns}
		\begin{definition}{\insertsubsection\mysource{\lientzswanson}}
			\mycite{Software maintenance performed to improve the performance, maintainability, or other attributes of a computer program.} \hfill \deutsch{perfektive Wartung}
		\end{definition}
		\begin{example}{Examples}
			\begin{itemize}
				\item better handling of large files in a text editor
				\item more efficient memory layout (e.g., in Linux)
				\item reduction of smells by means of refactorings
			\end{itemize}
		\end{example}
		\uncover<2->{\begin{example}{Meltdown/Spectre}
				\begin{itemize}
					\item January 2018 discovered
					\item caused by branch prediction in microprocessors
					\item fixed by software patches
					\item performance loss of up-to 30\%
				\end{itemize}
			\end{example}}
		\nextcolumn
		\uncover<3->{\begin{example}{CrowdStrike Outage\mysource{\href{https://en.wikipedia.org/wiki/CrowdStrike}{wikipedia.org}, \href{https://www.bbc.com/news/articles/c23k4yyjxp3o}{bbc.com}}}
				\begin{itemize}
					\item CrowdStrike's Falcon received update to improve \emph{security}
					\item update on July 19, 2024 affected 8.5 million devices running Windows 10 or 11
					\item systems were indeed more secure ;-)
					\item follow-up corrective maintenance was required including manual reboot
				\end{itemize}
				\picDark[width=\linewidth]{failures/crowdstrike} % TODO better identifying picture for crowdstrike?
			\end{example}}
	\end{fancycolumns}
\end{frame}
% TODO it is annoying that perfective maintenance has an overlap with preventive maintenance, as better maintainability typically also results in fewer problems

% TODO software aging (see Parnas-SoftwareAging, GodfreyGerman-LehmansLaws)

\begin{frame}{Adaptive Maintenance}
	\begin{fancycolumns}
		\begin{definition}{\insertsubsection\mysource{\lientzswanson}}
			\mycite{Software maintenance performed to make a computer program usable in a changed environment.} \hfill \deutsch{adaptive Wartung}
		\end{definition}
		\begin{example}{}
			desktop application for a new version of an operating system (e.g., from Windows 10 to 11)
		\end{example}
		\nextcolumn
		\picDark[width=\linewidth]{changes/windows7-11}
	\end{fancycolumns}
\end{frame}

\widexkcdframe{2038} % y2k bug vs 2038 bug

\begin{frame}{Y2K Problem \mytitlesource{\href{https://en.wikipedia.org/wiki/Year_2000_problem}{wikipedia.org}}}
	\begin{fancycolumns}[b]
		\begin{definition}{Y2K Problem}
			\begin{itemize}
				\item years were often stored by only two digits
				\item high awareness in media
				\item about 500 billion dollars spent
				\item about 90\% before 2000
				\item most major problems could be prevented
			\end{itemize}
		\end{definition}
		\centering\picDark[width=.8\linewidth]{failures/y2k-javascript}
		\nextcolumn
		\vspace{-15mm}

		\uncover<3->{\centering\pic[width=.5\linewidth]{failures/y2k}}
		\begin{example}{Selected Y2K Failures}
			\begin{itemize}
				\item<4-> 1999-01-01: computers for temporary passports in Sweden crashed
				\item<5-> 1999-12-28: credit and debit card transactions rejected because they could not be completed within four working days
					%\item 2000-01-01: date 2036-02-06 shown in nuclear power plant in Fukushima
				\item<2-> 2000-01-01: German bank accidentially transferred 6 million euro issued 1899-12-30
				\item<3-> 2000-01-03: french school showing wrong date
				\item<6-> 2020-01-01: New York City parking meters refuse credit cards
			\end{itemize}
		\end{example}
	\end{fancycolumns}
\end{frame}

\begin{frame}{Y2038 Problem \mytitlesource{\href{https://en.wikipedia.org/wiki/Year_2038_problem}{wikipedia.org}}}
	\begin{fancycolumns}
		\begin{definition}{Y2038 Problem}
			\begin{itemize}
				\item 2038-01-19 at 03:14:07 potential overflow
				\item systems affected storing Unix time as signed 32-bit integer
				\item similarly: 2106-02-07 for unsigned integers
				\item solution: use signed 64-bit integer
				\item 292 billion years to next overflow
			\end{itemize}
		\end{definition}
		\centering\picDark[width=\linewidth]{failures/y2038}
		\nextcolumn
		\begin{example}{Y2038 Failures}
			2006-05-13: AOL server software crashed
			\begin{itemize}
				\item it used 1 billion seconds extra to state that database transactions do not timeout
				\item \emph{temporary} solution: set the timeout to a lower value in configuration file
			\end{itemize}
			future failures? embedded systems without over-the-air updates are likely to be affected (e.g., ECUs in most automotive systems)
		\end{example}
	\end{fancycolumns}
\end{frame}

\begin{frame}{Preventive Maintenance}
	\begin{fancycolumns}
		\begin{definition}{\insertsubsection\mysource{\lientzswanson}}
			\mycite{Maintenance performed for the purpose of preventing problems before they occur.} \hfill \deutsch{präventive Wartung}
		\end{definition}
		\begin{example}{Y2K Problem}
			was very expensive, still not over
		\end{example}
		\begin{example}{Y2038 Problem}
			first failures already happen, will be expensive
		\end{example}
		\nextcolumn
		\begin{example}{Leap Years in Gregorian Calendar}
			\begin{itemize}
				\item solar year is 365.24 days
				\item every 4 years, but not every 100 years, but again every 400 years
				\item 2000 was a leap year
				\item one day extra: February 29, 2028
				\item common problem when durations are computed
				\item example: meat sold beginning of March
			\end{itemize}
		\end{example}
		\begin{example}{Leap Seconds}
			\begin{itemize}
				\item solar day longer than a day
				\item last leap second on 2016-12-31
				\item more details in SE1
			\end{itemize}
			\centering\pic[width=.5\linewidth,trim=0 0 0 0,clip]{changes/leap-second}
		\end{example}
	\end{fancycolumns}
\end{frame}

% Changelog (honest?)
% vgl. Release Management, Semantic Versioning
% Problems: time pressure, results in new failures, unpopular task

% TODO 22.4.2 Die Behandlung von Problemmeldungen \ludewiglichter

% TODO software erosion/rot/decay/entropy:  https://en.wikipedia.org/wiki/Software_rot

% TODO software estimation (bewertung)
% 	motivation, metrics for OOP
% 14 Metriken und Bewertungen \ludewiglichter

% New Content

% TODO APIs: every observable behavior will eventually be used by users

% TODO new example of Tesla: https://www.heise.de/news/Tesla-Autos-duerfen-nicht-mehr-furzen-Rueckruf-in-den-USA-6441936.html

% TODO add links to original literature: Chikofsky und Cross

\subsection{Reengineering}

\begin{frame}{\insertsubsection\ \mytitlesource{\ludewiglichter}}
	\begin{fancycolumns}
		\begin{definition}{Reengineering \mysource{Chikofsky und Cross}}
			\mycite{Reengineering is the examination and alteration of a subject system to reconstitute it in a new form and the subsequent implementation of the new form.}
		\end{definition}
		\begin{note}{Essence}
			combination of reverse engineering, refactoring, and forward engineering
		\end{note}
		\nextcolumn
		\begin{example}{Examples}
			\begin{itemize}
				\item reimplementing C programs in Rust
				\item modularizing a monolithic application into microservices
			\end{itemize}
		\end{example}
	\end{fancycolumns}
\end{frame}

\begin{frame}{Forward Engineering\ \mytitlesource{\ludewiglichter}}
	\begin{fancycolumns}[widths={60},animation=none]
		\begin{definition}{Forward Engineering \mysource{Chikofsky und Cross}}
			\mycite{Forward engineering is the traditional process of moving from high-level abstractions and logical, implementation-independent designs to the physical implementation of a system.}
		\end{definition}
		\uncover<2->{\begin{example}{Examples}
				\begin{itemize}
					\item modeling a system with activity diagrams or state machines \emph{based on} the requirements specification
					\item<3-> architecting a system with architectural patterns and component diagrams \emph{based on} the system specification
					\item<4-> designing a system with design patterns, class diagrams, sequence diagrams \emph{based on} the architecture specification
					\item<5-> implement a system in a particular programming language \emph{based on} the design specification
					\item<6-> \ldots
				\end{itemize}
			\end{example}}
		\nextcolumn
		\diagramForwardEngineering
	\end{fancycolumns}
\end{frame}

\begin{frame}{Reverse Engineering\ \mytitlesource{\ludewiglichter}}
	\begin{fancycolumns}[widths={40},animation=none]
		\diagramReverseEngineering
		\nextcolumn
		\begin{definition}{Reverse Engineering \mysource{Chikofsky und Cross}}
			\mycite{Reverse engineering is the process of analyzing a system to identify the system’s components and their interrelationships and create representations of the system in another form or at a higher level of abstraction.}
		\end{definition}
		\uncover<2->{\begin{example}{Examples}
				\begin{itemize}
					\item recovering which requirements have already been implemented
					\item<3-> exploring what classes are in the implementation and what is their purpose
					\item<4-> creating UML models for legacy software that was not modeled before
					\item<5-> adapting a UML model \emph{based on} changes in the implementation
					\item<6-> \ldots
				\end{itemize}
			\end{example}}
	\end{fancycolumns}
\end{frame}

\begin{frame}{Refactoring \mytitlesource{\ludewiglichter}}
	\vspace{-5mm}
	\begin{fancycolumns}
		\begin{definition}{Refactoring (Restructuring) \mysource{Chikofsky und Cross}}
			\mycite{Refactoring is a transformation from one form of representation to another at the same relative level of abstraction. The new representation is meant to preserve the semantics and external behavior of the original.}
		\end{definition}
		\begin{definition}{Refactoring \mysource{\refactoring}}
			\mycite{Refactoring is the process of changing a software system in such a way that it does not alter the external behavior of the code yet improves its internal structure. It is a disciplined way to clean up code that minimizes the chances of introducing bugs. In essence when you refactor you are improving the design of the code after it has been written.}
		\end{definition}
		\nextcolumn
		\begin{example}{Refactorings}
			\begin{itemize}
				\item rename X refactoring
				\item extract/pull-up method refactoring
			\end{itemize}
		\end{example}
		\begin{example}{Smells and Refactorings (Examples from last year SE2)}
			\pic[width=\linewidth]{2025st/wordcloud-smells-refactorings}
		\end{example}
	\end{fancycolumns}
\end{frame}

% TODO: styling
\begin{frame}
	\frametitle{Reengineering Exemplary Process \mytitlesource{\sommerville}}
	\tikzstyle{ioblock} = [rectangle, draw]
	\tikzstyle{processblock} = [rectangle, rounded corners, draw]
	\begin{figure}[htbp]
		\centering
		\begin{tikzpicture}[node distance=2cm and 0.8cm]
			\node[ioblock] (ori) at (0,0) {Original Software};
			\node[ioblock, above right = of ori] (documentation) {Documentation};
			\draw[->] (ori) -- node[sloped, above, midway] {\footnotesize Reverse Engineering} (documentation);
			\node[ioblock, below right = of ori] (refactored) {Refactored Software};
			\draw[->] (ori) -- node[sloped, above, midway] {\footnotesize Refactoring} (refactored);

			\node[processblock, below = of documentation] (modularization) {Modularization};
			\draw[->] (documentation) -- node[sloped, above, midway] {\footnotesize Forward Engineering} (modularization);
			\draw[->] (refactored) -- node[sloped, above, midway] {\footnotesize Forward Engineering} (modularization);

			\node[ioblock, right = 1.5cm of modularization] (updatedprogram) {Updated Program};
			\draw[->] (modularization) -- node[sloped, above, midway] {\footnotesize Forward Engineering} (updatedprogram);
			\node[ioblock, right = 1cm of updatedprogram] (newdata) {Updated Data};
			\node[ioblock, above = of newdata] (olddata) {Original Data};
			\draw[->] (olddata) -- node[sloped, above, midway] {\footnotesize Data Reengineering} (newdata);
			\draw[->] (updatedprogram) -- node[sloped, above, midway] {\footnotesize Data Reengineering} (newdata);
		\end{tikzpicture}

	\end{figure}
\end{frame}

\begin{frame}{Wisdom on Reengineering}
	\begin{fancycolumns}
		\pic[width=\linewidth,trim=50 0 0 0,clip]{people/kyle-simpson}
		\vspace{-7mm}

		\begin{note}{Kyle Simpson \mysource{\href{https://twitter.com/codewisdom/status/857988273674883072}{twitter.com}}}
			\mycite{There's nothing more permanent than a temporary hack.} \deutsch{Sprichwort: Nichts ist so beständig wie das Provisorium.} % TODO Quelle für deutsches Sprichwort suchen
		\end{note}
		\nextcolumn
		\vspace{11mm}
		\centering\pic[width=.5\linewidth,trim=0 0 0 0,clip]{people/jeff-sickel}
		\vspace{-3mm}

		\begin{note}{Jeff Sickel \mysource{\href{https://twitter.com/codewisdom/status/710046205317873664}{twitter.com}}}
			\mycite{Deleted code is debugged code.}
		\end{note}
	\end{fancycolumns}
\end{frame}

\subsection{Legacy Software}
\begin{frame}{\insertsubsection\ \deutsch{Altsysteme}}
	\begin{fancycolumns}
		\begin{definition}{Legacy Software \mysource{\ludewiglichter}}
			\mycite{A large software system that we don’t know how to cope with but that is vital to our organization.}
		\end{definition}
		\begin{exampletight}{Winamp 5.623 on 4k Display in Windows 10}
			\pic[width=\linewidth,trim=0 0 0 0,clip]{changes/winamp}
		\end{exampletight}
		\nextcolumn
		\begin{note}{Typical Properties \mysource{\ludewiglichter}}
			\begin{itemize}
				\item larger than 100k LOC
				\item older than 10 years
				\item original developers and architects not available anymore
				\item outdated programming languages and development concepts
				\item business-critical
				\item outdated or missing documentation
				\item based on outdated hardware and system software
				\item subject to numerous iterations of corrective and adaptive maintenance
				\item high costs for maintenance
			\end{itemize}
		\end{note}
	\end{fancycolumns}
\end{frame}

% \xkcdframe{2221} % emulated software, long-living

\subsection{Software Migration}
\begin{frame}{\insertsubsection\ \mytitlesource{\ludewiglichter}}
	\begin{fancycolumns}
		\begin{definition}{Software Migration}
			\begin{itemize}
				\item opposed to data and hardware migration
				\item renewing or replacing legacy software
				\item new software needs to be downward compatible
				\item new software meets current functional and non-functional requirements
				\item data often needs to be migrated
				\item outage as short as possible
			\end{itemize}
		\end{definition}
		\nextcolumn
		\begin{note}{Migration Strategies}
			\begin{itemize}
				\item big bang migration: whole legacy software is migrated in a single step
				\item incremental migration: stepwise renewal and replacement
				\item wrapping: building a new software around a stable version of the legacy software
				\item redevelopment: functionality implemented again and replaces legacy software
				\item combinations feasible (except for incremental and big bang)
			\end{itemize}
		\end{note}
	\end{fancycolumns}
\end{frame}

% TODO separate slides for migration strategies

\begin{frame}
	\frametitle{Big Bang Migration}
	\begin{fancycolumns}[widths={50,50}]
		\begin{definition}{Idea}
			\begin{itemize}
				\item Redevelop whole software
				\item Replace legacy software with new software in a single step
				\item No parallel operation of legacy and new software
			\end{itemize}
		\end{definition}
		\nextcolumn
		\begin{definition}{Advantages}
			\begin{itemize}
				\item Simple
				\item No overhead of maintaining two systems in parallel
			\end{itemize}
		\end{definition}
		\begin{note}{Disadvantages}
			\begin{itemize}
				\item Risk of failure
				\item Potential downtime
				\item Harder to identify and fix issues
			\end{itemize}
		\end{note}
	\end{fancycolumns}
\end{frame}

\begin{frame}
	\frametitle{Incremental Migration}
		\begin{fancycolumns}[widths={50,50}]
			\begin{definition}{Idea}
				\begin{itemize}
					\item Stepwise renewal and replacement of legacy software
					\item Parallel operation of legacy and new software
					\item Replace functionality on subsystem or module level 
				\end{itemize}
			\end{definition}
			\nextcolumn
			\begin{definition}{Advantages}
				\begin{itemize}
					\item Lower risk of failure
					\item Easier to identify and fix issues
					\item Not a single potentially long downtime
					\item Can already extend redeveloped parts with new features
				\end{itemize}
			\end{definition}
			\begin{note}{Disadvantages}
				\begin{itemize}
					\item More complex
					\item Higher overhead of maintaining two systems in parallel
					\item Requires careful planning and coordination
				\end{itemize}
			\end{note}
		\end{fancycolumns}
\end{frame}

\begin{frame}[fragile]
	\frametitle{Migration by Wrapping}
	\begin{fancycolumns}[widths={50,50}]
		\begin{definition}{Idea}
			\begin{itemize}
				\item Build a new software around a stable version of the legacy software
				\item New software provides an interface to the legacy software
				\item Call legacy functions from new software
			\end{itemize}
		\end{definition}
		\begin{example}{Example}
			\vspace{-3mm}
			\begin{lstlisting}[language=Java,basicstyle=\footnotesize]
// Invoke legacy operations A, B from new software
output = doSomeNewOperations(input);
invokeLegacyOperationA(convertToLegacyFormat(output));
invokeLegacyOperationB(convertToLegacyFormat(output));

// Update after Incremental migration
output = doSomeNewOperations(input);
doNewA(output);
invokeLegacyOperationB(convertToLegacyFormat(output));
			\end{lstlisting}
			\vspace{-3mm}
		\end{example}
		\nextcolumn
		\begin{definition}{Advantages}
			\begin{itemize}
				\item Lower reverse engineering effort
				\item Old functionality is preserved
				\item Can be performed step-wise
				\item New features can be added in the new software
			\end{itemize}
		\end{definition}
		\begin{note}{Disadvantages}
			\begin{itemize}
				\item Old systems are still running
				\item Maintenance of two systems in parallel
				\item Changes may still require editing the legacy software
			\end{itemize}
		\end{note}
	\end{fancycolumns}
\end{frame}


%\subsection{Big Bang Migration}
%
%\subsection{Migration by Wrapping}
%
%\subsection{Incremental Migration}
%
%\subsection{Redevelopment?}

% TODO reasons/requirements/process for migration

% TODO end of life: "the perfect race car crosses the finish line in first place and immediately falls into pieces" Ferdinand Porsche, https://www.fahrtraum.at/en/the-best-quotes-and-sayings-from-ferdinand-porsche/

\begin{frame}{End of Service}
	\begin{fancycolumns}
		\begin{exampletight}{{End of Life: December 31, 2020}}
			\centering\picDark[width=.85\linewidth,trim=0 0 0 0,clip]{changes/flash-eol}
		\end{exampletight}
		\nextcolumn
		\begin{exampletight}{Online Game Fuchstreff Discontinued}
			\picDark[width=\linewidth,trim=0 0 0 0,clip]{changes/fuchstreff}
		\end{exampletight}
	\end{fancycolumns}
\end{frame}

% TODO 32bit architectures discontinued in many Linux distributions
% https://www.debugpoint.com/32-bit-linux-distributions/

% TODO release, maintenance, end of life / alternatives to maintenance https://en.wikipedia.org/wiki/Software_maintenance

